%%
%% This is file `sigconf.tex',
%% generated with the docstrip utility.
%%
%% The original source files were:
%%
%% samples.dtx  (with options: `all,proceedings,sigconf')
%% 
%% IMPORTANT NOTICE:
%% 
%% For the copyright see the source file.
%% 
%% Any modified versions of this file must be renamed
%% with new filenames distinct from sigconf.tex.
%% 
%% For distribution of the original source see the terms
%% for copying and modification in the file samples.dtx.
%% 
%% This generated file may be distributed as long as the
%% original source files, as listed above, are part of the
%% same distribution. (The sources need not necessarily be
%% in the same archive or directory.)
%%
%%
%% Commands for TeXCount
%TC:macro \cite [option:text,text]
%TC:macro \citep [option:text,text]
%TC:macro \citet [option:text,text]
%TC:envir table 0 1
%TC:envir table* 0 1
%TC:envir tabular [ignore] word
%TC:envir displaymath 0 word
%TC:envir math 0 word
%TC:envir comment 0 0
%%
%% The first command in your LaTeX source must be the \documentclass
%% command.
%%
%% For submission and review of your manuscript please change the
%% command to \documentclass[manuscript, screen, review]{acmart}.
%%
%% When submitting camera ready or to TAPS, please change the command
%% to \documentclass[sigconf]{acmart} or whichever template is required
%% for your publication.
%%
%%
\documentclass[sigconf]{acmart}
%%
%% \BibTeX command to typeset BibTeX logo in the docs
\AtBeginDocument{%
  \providecommand\BibTeX{{%
    Bib\TeX}}}

%% Rights management information.  This information is sent to you
%% when you complete the rights form.  These commands have SAMPLE
%% values in them; it is your responsibility as an author to replace
%% the commands and values with those provided to you when you
%% complete the rights form.
\setcopyright{acmlicensed}
\copyrightyear{2018}
\acmYear{2018}
\acmDOI{XXXXXXX.XXXXXXX}
%% These commands are for a PROCEEDINGS abstract or paper.
\acmConference[Conference acronym 'XX]{Make sure to enter the correct
  conference title from your rights confirmation email}{June 03--05,
  2018}{Woodstock, NY}
%%
%%  Uncomment \acmBooktitle if the title of the proceedings is different
%%  from ``Proceedings of ...''!
%%
%%\acmBooktitle{Woodstock '18: ACM Symposium on Neural Gaze Detection,
%%  June 03--05, 2018, Woodstock, NY}
\acmISBN{978-1-4503-XXXX-X/2018/06}


%%
%% Submission ID.
%% Use this when submitting an article to a sponsored event. You'll
%% receive a unique submission ID from the organizers
%% of the event, and this ID should be used as the parameter to this command.
%%\acmSubmissionID{123-A56-BU3}

%%
%% For managing citations, it is recommended to use bibliography
%% files in BibTeX format.
%%
%% You can then either use BibTeX with the ACM-Reference-Format style,
%% or BibLaTeX with the acmnumeric or acmauthoryear sytles, that include
%% support for advanced citation of software artefact from the
%% biblatex-software package, also separately available on CTAN.
%%
%% Look at the sample-*-biblatex.tex files for templates showcasing
%% the biblatex styles.
%%
%%
%% The majority of ACM publications use numbered citations and
%% references.  The command \citestyle{authoryear} switches to the
%% "author year" style.
%%
%% If you are preparing content for an event
%% sponsored by ACM SIGGRAPH, you must use the "author year" style of
%% citations and references.
%% Uncommenting
%% the next command will enable that style.
%%\citestyle{acmauthoryear}


%%
%% end of the preamble, start of the body of the document source.
\begin{document}

%%
%% The "title" command has an optional parameter,
%% allowing the author to define a "short title" to be used in page headers.
\title{Test-Time Intervention Cascades for Heterogenous Workloads}

%%
%% The "author" command and its associated commands are used to define
%% the authors and their affiliations.
%% Of note is the shared affiliation of the first two authors, and the
%% "authornote" and "authornotemark" commands
%% used to denote shared contribution to the research.
\author{David Warren}
\email{david.warren@hdr.mq.edu.au}
\orcid{\TODO}
\author{Mark Dras}
\correspondingauthor
\email{mark.dras@mq.edu.au}
\affiliation{%
  \institution{Macquarie University}
  \city{Sydney}
  \state{NSW}
  \country{Australia}
}




%%
%% By default, the full list of authors will be used in the page
%% headers. Often, this list is too long, and will overlap
%% other information printed in the page headers. This command allows
%% the author to define a more concise list
%% of authors' names for this purpose.
\renewcommand{\shortauthors}{Warren \& Dras}

%%
%% The abstract is a short summary of the work to be presented in the
%% article.
\begin{abstract}
  
\end{abstract}

%%
%% The code below is generated by the tool at http://dl.acm.org/ccs.cfm.
%% Please copy and paste the code instead of the example below.
%%
\begin{CCSXML}
\TODO
<ccs2012>
 <concept>
  <concept_id>00000000.0000000.0000000</concept_id>
  <concept_desc>Do Not Use This Code, Generate the Correct Terms for Your Paper</concept_desc>
  <concept_significance>500</concept_significance>
 </concept>
 <concept>
  <concept_id>00000000.00000000.00000000</concept_id>
  <concept_desc>Do Not Use This Code, Generate the Correct Terms for Your Paper</concept_desc>
  <concept_significance>300</concept_significance>
 </concept>
 <concept>
  <concept_id>00000000.00000000.00000000</concept_id>
  <concept_desc>Do Not Use This Code, Generate the Correct Terms for Your Paper</concept_desc>
  <concept_significance>100</concept_significance>
 </concept>
 <concept>
  <concept_id>00000000.00000000.00000000</concept_id>
  <concept_desc>Do Not Use This Code, Generate the Correct Terms for Your Paper</concept_desc>
  <concept_significance>100</concept_significance>
 </concept>
</ccs2012>
\end{CCSXML}

\TODO
\ccsdesc[500]{Do Not Use This Code~Generate the Correct Terms for Your Paper}
\ccsdesc[300]{Do Not Use This Code~Generate the Correct Terms for Your Paper}
\ccsdesc{Do Not Use This Code~Generate the Correct Terms for Your Paper}
\ccsdesc[100]{Do Not Use This Code~Generate the Correct Terms for Your Paper}

%%
%% Keywords. The author(s) should pick words that accurately describe
%% the work being presented. Separate the keywords with commas.
\TODO
\keywords{Do, Not, Use, This, Code, Put, the, Correct, Terms, for,
  Your, Paper}
%% A "teaser" image appears between the author and affiliation
%% information and the body of the document, and typically spans the
%% page.

\TODO
\received{20 February 2007}
\received[revised]{12 March 2009}
\received[accepted]{5 June 2009}

%%
%% This command processes the author and affiliation and title
%% information and builds the first part of the formatted document.
\maketitle


\section{Introduction}

Modern Large Language Model (LLM) deployments operate under challenging multi-objective performance constraints with complex and unpredictable inputs. LLMs are required to provide high-quality, factually accurate, instruction-compliant, and appropriately formatted outputs while minimizing wasteful computational overhead, latency, and cost \cite{PLACEHOLDER:llm-deployment-constraints}. These operational requirements have driven significant research into inference efficiency through adaptive computation, which can improve the cost-performance frontier without requiring model training or fine-tuning at the LLM scale \cite{efficiencysurvey, moslem2026dynamicmodelroutingcascading}. Two related adaptive LLM inference paradigms have emerged in the literature: model cascading \cite{varshney-baral-2022-model} and model routing \cite{ding2024hybrid, ong2024routellmlearningroutellms, routerbench}. Model cascades utilize an inference-first sequential escalation approach to process samples with small models where appropriate, reserving expensive computation for difficult samples, whereas model routing leverages information in the input space to direct samples to appropriately capable models prior to primary inference.
 
Traditional model cascades provide advantageous cost-performance trade-offs in Natural Language Processing (NLP) classification and some generation tasks \cite{varshney-baral-2022-model, Gupta2024}, but inherit critical architectural constraints from their classical machine learning origins: a single, post-inference escalation action representing improved general competence. Classical cascading frameworks equate capability improvement strictly with parameter-scale escalation, moving monotonically from a cheap, weak model to a more expensive, more capable one whenever the cheap model's output is judged insufficient. This approach becomes less efficient at lower inter-tier cost ratios due to the unconditional execution of the first model, including on samples subsequently deferred. As additional stages are inserted within a fixed cost range, earlier execution costs accumulate while the incremental benefit of each stage may diminish. This formulation is well suited to more homogeneous, traditional machine learning contexts, in which a model's failures are reasonably well predicted by its position on a single, sparsely populated capacity axis. However, this approach is a poor match for LLMs, whose capability deficits are not well summarized by a single scalar notion of "more capable" and many diverse escalations are available.
 
While this mismatch limits the utility of cascades naively applied, LLMs have unique properties with no clear analogue in the classical machine learning that enable significant departures from the traditional cascade paradigm. In-context learning (ICL) allows an LLM's behavior to be substantially altered at inference time, without any parameter update, purely through the content of its input context \cite{icl}. This is the mechanism underlying a wide array established context-based, non-parametric test-time interventions in LLMs. Crucially, these interventions do not retrain or fine-tune the underlying model; all of them rely on the same underlying capacity of the model to be steered by its context. Cascades of traditional machine learning models have no equivalent to this; a boosted decision tree or a convolutional network cascade cannot be "prompted" into recovering a different subset of errors without retraining a new stage. The availability of ICL therefore introduces additional axes of capability improvement, requiring no additional parameters or training, and typically available at a cost far below that of escalating to a larger model. This fundamental shift in failure modes and escalation options motivates our departure from traditional cascade design.
 
There are many ways to potentially partition the space of LLM capability deficits; we do not propose or rely on a canonical taxonomy here. As illustrative motivation, some failures stem from missing factual knowledge, others from incorrect reasoning steps despite adequate knowledge, and others from an inability to comply with formatting or structural constraints despite understanding the task. Modern LLM system design has produced non-parametric test-time interventions that loosely target failures of these kinds without escalating model scale or requiring any fine-tuning: Retrieval-Augmented Generation (RAG) is commonly used to address factual knowledge gaps \cite{NEURIPS2020_6b493230}; self-verification and chain-of-thought-style prompting are commonly used to address reasoning errors \cite{selfrefine, selfverify}; and few-shot prompt expansion is commonly used to address structural or formatting non-compliance \cite{icl}. We select these three interventions as examples of comparably-priced, non-parametric actions with distinct typical use cases, not as an exhaustive or mutually exclusive set — in practice, their recovery sets overlap substantially rather than partitioning cleanly (self-verification and few-shot expansion, for instance, often recover overlapping structural errors). What matters for the remainder of this work is only that their execution costs are comparable to one another and substantially lower than escalating to a larger model tier, not that they map onto any particular capability taxonomy.
 
\subsection{From sequential stacking to routing and cascading}
 
This landscape of diverse, comparably-priced interventions presents an opportunity for efficient cascades, but also a design challenge. Placing every intervention in a separate sequential stage incurs the cost of earlier actions before a sample can access later ones. Interventions have situational benefits and overlapping recovery sets and expanding the cascade to include these interventions may not expand the recoverable error set proportionately. The fundamental conflict is then: additional choices are valuable, but requiring their sequential execution may consume that value. \citet{jitkrittum} describe a related challenge in the two-model setting when the downstream model behaves as a specialist. We examine this challenge across several similarly priced interventions, testing both their complementarity and the marginal benefit of executing additional stages.

 When actions incur similar costs but resolve distinct, overlapping failures, selection among them can be placed within a single routed tier, extracting the additional value where possible while constraining the cost to a single inference. This motivates our central architectural claim: routing and cascading are not competing solutions to the same allocation problem, but compatible mechanisms suited to two different decisions. Routing determines which intervention to execute; cascading determines whether to continue from the current response. Existing hybrid systems already combine these mechanisms in some fashion: AutoMix uses answer verification to inform model escalation \cite{automix}, while cascade routing selects among candidate model sequences and can bypass models in a conventional ordered cascade \cite{dekoninck2025unified}. Our objective is different from making the sequence of model calls more flexible. We introduce a tier interface that makes multiple similarly priced test-time interventions available as alternatives, charging a sample for only the selected intervention rather than requiring it to traverse the others. This allows intervention diversity to expand without a corresponding increase in cascade depth, while preserving a separate escalation decision across the substantially larger cost gap to the large model.
 
This division of labor follows from a difference in the mechanics of the two decisions, not merely a difference in vocabulary. Cascading is fundamentally a post-inference decision: the small model is run to completion, and a calibrated confidence or uncertainty signal derived from its own output is used to decide whether to defer to a larger model \cite{varshney-baral-2022-model}. Its effectiveness therefore depends on the specific performance and cost relationship between the two tiers being bridged: on how well the small model's confidence is calibrated to its actual correctness, and on how large the accuracy gain from escalation is relative to its cost at the operating point chosen \cite{varshney-baral-2022-model, PLACEHOLDER:gupta-chow-cascades}. This dependence is precisely why confidence-based deferral can fail in the settings identified by \citet{jitkrittum2023when}, where the escalation target behaves as a specialist rather than a uniformly better generalist, or where confidence is miscalibrated under distribution shift. Routing, by contrast, is fundamentally a pre-inference decision: a router assigns a query to a computation path using only the query itself (and, in our mid-cascade instantiation, the small model's output) before any candidate downstream output has been produced \cite{chen2023frugalgptuselargelanguage, routerbench}. Its effectiveness therefore depends primarily on how accurately the router triages the incoming query distribution across the available options, rather than on a calibrated confidence-accuracy relationship between any specific pair of options; because the routing decision does not require observing an intermediate output from each candidate path, it is comparatively insensitive to how many paths are available or how they are ordered, in contrast to a cascade's sequential dependence on a specific tier ordering \cite{chen2023frugalgptuselargelanguage, routerbench}. It is this difference in mechanics that motivates treating routing and cascading as complementary rather than substitutable: the intervention tier, where several comparably-priced and imperfectly-substitutable actions are available, is well suited to triage-based routing, while the large disparity between the intervention tier and the large model is well suited to calibrated post-inference deferral.
 
We instantiate this composition as a hybrid architecture with two components situated on either side of the intervention tier. The first is an intra-tier router, which selects among comparably-priced, heterogeneous test-time interventions at the sample level. Rather than routing purely ex-ante from the prompt alone, our router is positioned as a hybrid mid-cascade observer: a lightweight learned classifier that takes as input both the raw query and the small model's initial output, and learns the empirical relationship between that query-output pair and which intervention is likely to recover the error, without observing any intervention's outcome before committing to one. The second is an inter-tier deferral mechanism, which determines after the selected intervention has been applied whether the sample should still be escalated to the substantially more expensive large model. We deliberately reuse a conventional, non-learned uncertainty-based deferral rule for this decision, applied identically both to a conventional small-to-large model cascade baseline and to the post-intervention escalation decision in our system, so that any improvement we observe is attributable to the routing and intervention layer rather than to a different deferral policy across conditions. The two decisions are evaluated independently throughout this work, since conflating "did the router select a useful intervention" with "was escalation to the large model necessary" would obscure which mechanism is responsible for any observed gain.
 
\subsection{Evaluation methodology}
 
Evaluating a routing mechanism on a single, narrow benchmark understates the situational variability that routing is designed to exploit: if a workload tests one deliberately isolated capacity, one intervention is likely to dominate throughout, and the routing decision degenerates into a fixed choice. We therefore evaluate on two workloads that are themselves mixtures of underlying tasks, rather than single-capacity benchmarks, so that the value of situational selection can actually be observed. The first, MixEval \cite{mixeval}, is a ground-truth-gradable benchmark constructed to approximate a naturalistic, real-world query distribution by matching web-mined user queries against existing benchmark items. The second is a pre-mixed, diverse collection of tasks drawn from RouterBench \cite{routerbench}, which spans commonsense reasoning, knowledge-intensive question answering, mathematical reasoning, and code generation rather than a single task family. The two workloads differ in the structure of their heterogeneity — a naturalistic, blended query distribution versus an aggregation of distinct benchmark suites — which additionally lets us test whether the value of routing scales with workload heterogeneity.
 
For cost accounting, we adopt a deliberately simple model based on separate per-token rates for prompt and completion tokens, in the manner used by commercial LLM API providers and by prior LLM routing and cascading work \cite{chen2023frugalgptuselargelanguage, routerbench}, rather than a multi-dimensional simulation of memory, compute, and latency. Cost accounting is inference focused and comparable across interventions with otherwise different cost profiles (e.g., expanded context versus multi-pass generation), at the expense of abstracting away lower-level systems effects such as batching or memory pressure that a full hardware simulation would capture.
 
\subsection{Contributions}
 
This work makes three central contributions:
\begin{enumerate}
\item \textbf{Situational intervention complementarity}: individual test-time interventions provide efficient but non-uniform, partially overlapping recovery sets relative to the small model's errors, such that no single intervention is universally optimal. We additionally compare these recovery sets against the errors recovered by escalating to the large model, to assess whether test-time interventions recover a meaningfully different set of failures than parameter-scale escalation does, rather than a cheaper subset of the same failures.
\item \textbf{The inefficiency of sequential stacking}: naively arranging these interventions into a linear escalation pipeline, in the manner of a traditional cascade, produces rapidly diminishing returns and unnecessary computational cost, motivating a different allocation mechanism within the intervention tier.
\item \textbf{Pareto superiority of the routed cascade}: our combined architecture, using intra-tier routing among interventions together with conventional inter-tier deferral to the large model, yields a significant Pareto improvement over both the conventional small-to-large model cascade and the linear intervention-stacking ablation, on both evaluated workloads.
\end{enumerate}
 
Together, these contributions extend LLM cascades beyond a single, model-scale escalation axis. Instead, different computation-allocation mechanisms are appropriate at different points in a structured cost hierarchy: sample-level routing where actions are comparably priced and situational, and post-inference cascading where the available escalation crosses a substantial cost disparity.

\section{Related Work}
\label{sec:related-work}

Our work connects two cost-optimizing adaptive inference paradigms, cascading and routing, and the diverse non-parametric test-time interventions available to LLMs. This depends heavily on how mechanisms uniquely applicable to LLMs interact with adaptive inference and the opportunities and constraints that emerge as a result. We discuss here the adaptive LLM inference literature and briefly cover a limited, illustrative set of test-time interventions for LLMs. These interventions are selected because of their established value in research and industry, accessibility and challenging but typical cost profiles with respect to cascading, not because of their optimality or comprehensive coverage.

\subsection{Model cascades and confidence-based deferral}

Cascades allocate computation sequentially, allowing inexpensive stages to handle some inputs while reserving further processing for others. The decision to defer to a more expensive model occurs after, and is informed by, the inference at the current stage. This principle predates language models and has been used to improve inference efficiency in a range of domains \citet{viola, todo}. In NLP, \citet{varshney-baral-2022-model} evaluated model cascading as a means of improving efficiency and accuracy in a wide range of configurations, establishing the viability of the paradigm applied to language models. Later work

%Later work FrugalGPT learns cascades over language-model APIs with different costs and capabilities \cite{chen2024frugalgpt}. These approaches establish the value of observing an inexpensive prediction before paying for another model.

The effectiveness of this strategy depends on the relationship between current and downstream errors as well as the costs incurred at each step. \citet{jitkrittum} characterize optimal cascade deferral and identify settings in which current-model confidence is insufficient, including specialist downstream models, label noise, and distribution shift. Their analysis is particularly relevant when an escalation helps only a subset of inputs: uncertainty about the current answer does not identify the expected benefit of the available successor. For generative language models, \citet{Gupta2024} examine deferral using sequence and token-level uncertainty, showing the importance of how uncertainty is aggregated across generated tokens. Together, these studies motivate separating an answer-quality signal from a prediction of downstream usefulness.

A further constraint is the cost of sequential access. The first stage is executed unconditionally, including for queries that will ultimately be answered downstream. Sequential cascaded execution therefore becomes less efficient as the inter-tier cost ratio narrows, and adding stages within a fixed cost range requires sufficient conditional gains to offset accumulated execution costs. This is a practical limitation rather than a requirement that cascade stages differ in parameter count. It motivates grouping similarly priced options behind a selector instead of assigning each its own sequential tier.

Our focus is the downstream action set. We consider several context-based interventions applied to the same underlying model, whose potential benefits need not admit a common ordering. The question is whether selecting among these actions can improve the quality--cost trade-off before a subsequent escalation to a larger model. Conventional uncertainty-based deferral supplies a controlled component of this comparison; the contribution concerns intervention complementarity and allocation rather than a new confidence estimator.

\subsection{Routing and hybrid model selection}

Query routing chooses a model without first obtaining outputs from candidates, using the relationship between output quality and the query itself. RouteLLM learns routing policies from preference data to balance response quality and inference cost \cite{ong2024routellmlearningroutellms}. RouterBench provides an evaluation framework and a collection of recorded model outcomes for comparing routing strategies across heterogeneous tasks \cite{routerbench}. These formulations make model diversity useful when relative performance varies across inputs, rather than assuming one model is preferable throughout a workload. The strength of this approach relies 

Hybrid allocation methods have different objectives and constraints. AutoMix uses an answer and a few-shot verification signal to inform model escalation \cite{automix}. \citet{dekoninck2025unified} define a supermodel as a sequence of model calls and formulate cascade routing as selection among supermodels consistent with calls already made. This permits bypassing intermediate models and reconsidering which model to execute next, rather than always taking the next step in a fixed cascade. Their formulation optimizes model-sequence selection under a quality--cost objective, accounting for alternative combinations of calls at each decision point.

Our focus is the inclusion of multiple similarly priced interventions as a single routed tier. Rather than adding each option as another sequential stage, the tier executes one selected intervention and then exposes its answer to a conventional deferral rule. Its purpose is to make a wider suite of test-time interventions economically accessible within a cascade without requiring a wider range of intervention costs. The distinction is therefore the tier construction and its empirical motivation, rather than the general possibility of combining routing and cascading. We evaluate the overlap of intervention recovery sets, the cost of sequential access to those alternatives, and the benefit of selecting one before a separate large-model escalation.

\subsection{Test-time interventions and adaptive computation}

In-context learning demonstrates that task behaviour can be changed through examples without updating model parameters \cite{icl}. Retrieval-augmented generation supplies external evidence for knowledge-intensive tasks \cite{NEURIPS2020_6b493230}, while chain-of-thought prompting elicits intermediate reasoning through demonstrations \cite{wei2022chain}. Output-conditioned interventions offer another form of adaptation: Self-Refine iterates between feedback and revision with the same model \cite{selfrefine}, and self-verification uses backward checks to assess candidate reasoning answers \cite{selfverify}. These methods establish useful intervention mechanisms, but their average benefits do not imply that they improve every answer or address disjoint failure categories. In particular, \citet{huang2024selfcorrect} find that intrinsic self-correction can fail or degrade performance on the reasoning settings they investigate. We therefore treat recovery, overlap, and regressions as empirical properties rather than assigning each intervention an exclusive capability label.

Adaptive retrieval is especially close to our allocation problem. Adaptive-RAG uses a classifier to select among no retrieval, single-step retrieval, and iterative retrieval according to question complexity \cite{jeong2024adaptiverag}. Self-RAG trains a language model to control retrieval and assess generated content using reflection tokens \cite{asai2024selfrag}. Both recognize that retrieval should be conditional on the input and generation state. Our intervention tier broadens the alternatives beyond retrieval strategies to include verification and few-shot expansion, and places this choice inside a cascade with a separate large-model fallback. The generation model remains fixed when applying these interventions, although the intervention selector is learned.

More broadly, test-time scaling investigates how additional inference computation can substitute for increased model capacity. \citet{snell2025testtime} study adaptive allocation of revision and search computation for reasoning, finding that the effective allocation depends on problem difficulty. BEST-Route jointly selects a model and the number of responses to sample according to query difficulty and quality requirements \cite{ding2025bestroute}, extending routing to include a sampling budget. These studies motivate evaluating computation policies at matched budgets rather than comparing only their unconstrained accuracy. Our question concerns the composition of distinct intervention types within a mixed workload: whether selective execution offers a better trade-off than applying them in a fixed sequence. Neither universal intervention superiority nor universally harmful sequential refinement is assumed.

\subsection{Heterogeneous Workloads}

Workload composition determines whether conditional selection has room to improve over a fixed action. RouterBench explicitly supports comparisons across diverse tasks and model outcomes \cite{routerbench}, while MixEval constructs benchmark mixtures intended to approximate real-world evaluation demands \cite{mixeval}. These resources motivate evaluating intervention selection on mixtures rather than relying exclusively on a benchmark that favours one intervention. They also make aggregate performance insufficient for explaining an allocation policy: similar overall scores can conceal different recovered and newly introduced errors.

Accordingly, our study connects three analyses: the complementarity of intervention outcomes, the additional cost of sequential intervention execution, and the quality-cost frontier of a routed intervention tier with subsequent model escalation. This positioning builds on established hybrid allocation methods while examining the specific value of heterogeneous, non-parametric interventions as intermediate cascade actions.

\section{Cascade Escalation Efficiency Concepts}


\section{Methods}
\label{sec:methods}

We study whether a routed intervention tier can improve a small-to-large model cascade by making several similarly priced test-time actions available without executing them in sequence. The experimental design separates three questions: whether interventions recover complementary failures, whether sequential execution uses that complementarity efficiently, and whether a learned selector converts it into a better quality-cost frontier. 
\subsection{Problem formulation and tier structure}

Let $\mathcal{D}=\{(x_i,z_i)\}_{i=1}^{N}$ be a workload of queries and reference answers or task-specific evaluation targets. A small model $M_S$ first produces $y_{i,0}=M_S(x_i)$. The intervention set is $\mathcal{A}=\{\mathrm{rag},\mathrm{verify},\mathrm{fewshot}\}$, with response $y_{i,a}=M_S(T_a(x_i,y_{i,0}))$ for intervention-specific context transformation $T_a$. Retrieval and few-shot expansion use the query, while verification also uses the initial answer. A separate large model $M_L$ provides the terminal response $y_{i,L}=M_L(x_i)$. Let $q_{i,a}\in\{0,1\}$ denote the correctness of a recorded response under the workload's scoring rule.

The system has three logical tiers: initial small-model inference, a routed intervention tier, and large-model fallback. An uncertainty gate first either accepts $y_{i,0}$ or admits the query to the intervention tier. For admitted queries, the selector chooses exactly one action $a_i=\pi_\theta(x_i,y_{i,0})$. A second uncertainty gate accepts the resulting intervention response or defers to $M_L$. Thus a deployed path executes zero or one intervention; increasing the number of available interventions does not require traversing additional intervention stages. The selector uses neither reference answers nor candidate intervention outputs at inference time.

This tier structure addresses the economics of sequential escalation. Sequential cascaded execution becomes less efficient at lower inter-tier cost ratios because the first model is executed unconditionally, including on samples ultimately processed downstream. Adding similarly priced sequential stages increases the accumulated cost without necessarily expanding the overall cost range or recovering many additional errors. The routed tier makes these alternatives accessible through selection while retaining uncertainty-based deferral across the larger cost gap to $M_L$. This motivates the architecture without assuming that a particular cost ratio alone determines whether cascading is worthwhile.

\subsection{Models and intervention instantiations}

The experimental configuration uses Qwen3-8B as $M_S$ and Llama-3.3-70B-Instruct as $M_L$. The generation model's parameters remain fixed across interventions. Non-thinking runs use greedy decoding without beam search; prompt templates and generation limits are task dependent. In the mixed RouterBench configuration, code and mathematical tasks permit up to 512 generated tokens, while other task families use a 128-token limit. These are configured limits, not assumptions about realized response lengths.


\paragraph{Retrieval augmentation.}
The retrieval implementation uses a hybrid text retrieval strategy over a BM25 index and a FAISS dense index using \texttt{all-MiniLM-L6-v2} embeddings. Results from the two retrieval channels are combined by reciprocal-rank fusion with rank constant 60 \cite{todo}. The generation notebook requests five passages for multiple-choice or mixed workloads, retaining passages with fused score at least 0.015. Retained text is prepended to the query with instructions to answer directly and to disregard irrelevant evidence. The retrieval corpus is a limited prebuilt Wikipedia index representing broad general knowledge \cite{todo}.

\paragraph{Verification.}
The verification action makes an additional call to $M_S$ containing the original query and saved baseline answer \cite{selfverify}. Its prompt asks the model to reproduce a correct answer or replace an incorrect or incomplete answer, returning only the final answer. It does not provide an external correctness label. Consequently, verification may preserve, repair, or damage the initial answer, all of which are retained in the outcome analysis.

\paragraph{Few-shot expansion.}
The few-shot action prepends five question-answer demonstrations before regenerating an answer with $M_S$. Demonstrations are drawn from samples outside the evaluable dataset and are fixed for the entire workload rather than tailored to the task relevant to the instance.

The three actions instantiate a broader tier interface rather than an exhaustive intervention set. Inclusion requires an executable action, an observable response and cost, and outcomes used for training the selector; it does not require assigning each action an exclusive error category.

\subsection{Outcome collection and workload scoring}

The generation process records prompts, responses, token sequences, retained generation scores, reference targets, and task metadata for each configuration. The merging stage derives response-quality and uncertainty columns and assembles the outcome tables used by the cascade simulator. Comparisons use the same queries across baseline, intervention, and large-model conditions, so that structural and policy differences can be fairly evaluated on paired outcomes. 

We consider MixEval's multiple-choice and free-form queries \cite{mixeval} and a heterogeneous RouterBench-derived workload \cite{routerbench}. The Routerbench workload is a subset of the published Routerbench dataset that includes diverse automatically evaluable tasks drawn form: MMLU, ARC-Challenge, HellaSwag, WinoGrande, grade-school mathematics, and MBPP code-generation\cite{todo}. The processing code evaluates multiple-choice answers against normalized answer options, mathematical answers through numeric-answer extraction, and MBPP responses against the provided tests. Samples with free-form answers are evaluated using GPT5-mini as a judge with reference answers to allow for flexible expression and alternate wording of correct answers \cite{todo gpt judge}. 

While exhaustive inference over all models and interventions on all samples samples is used for the sake of constructing the experimental data, cascade evaluation reflects the adaptive inference cost under a given deferral rate and policy. The cascade cost at each rate considers only the inference that a policy would have executed at the specific operating point and accumulates the costs of those actions. These costs include the inference stages that are conducted but not scored under a cascaded inference t does not refund earlier calls after escalation. The terminal response is the answer produced at the stopping stage, without reference-assisted selection among earlier answers.

\subsection{Learning the intra-tier selector}

The selector is a DeBERTa-v3-small model receiving the query concatenated with the initial response. Attention-masked mean pooling, dropout of 0.1, and a linear head produce one logit per intervention. Independent sigmoid outputs allow several interventions to be useful on the same query. This matches the partially overlapping recovery sets rather than forcing a unique intervention label for every training example.

In the three-way configuration, the training target for action $a$ is its binary destination correctness, $t_{i,a}=q_{i,a}$. This differs from labelling only strict improvements over the baseline: an intervention that preserves a correct answer also has a positive target. Training retains informative examples on which the intervention labels disagree,
\begin{equation}
\mathcal{I}=\{i:\min_{a\in\mathcal{A}}t_{i,a}<\max_{a\in\mathcal{A}}t_{i,a}\},
\end{equation}
and minimizes an unweighted binary cross-entropy objective,
\begin{equation}
\mathcal{L}(\theta)=\frac{1}{|\mathcal{I}|\,|\mathcal{A}|}
\sum_{i\in\mathcal{I}}\sum_{a\in\mathcal{A}}
\operatorname{BCE}(t_{i,a},s_{i,a}),
\end{equation}
where $s_{i,a}=\sigma(f_\theta(x_i,y_{i,0})_a)$. At inference, $\pi_\theta(x_i,y_{i,0})=\arg\max_a s_{i,a}$, with ties resolved by the fixed action order. These scores rank actions; because training is restricted to informative examples, we do not assume they are calibrated correctness probabilities for the full workload.

TODO: Move these params to appendix once we have one.

A 5-fold cross-validation strategy is used for training to provide out-of-fold predictions used in evaluation. For each fold, the selector is fitted on informative examples from the other four folds and evaluated on every example in the held-out fold, including examples with identical intervention labels. Fold assignments are shared across model and intervention. This permits evaluation using predictions from a selector that was not fitted on the evaluated query.


\subsection{Inter-tier deferral and policy replay}

The primary deferral signal is an uncertainty score derived from the mean log probability of generated tokens \cite{Gupta2024}. For an answer containing $T$ generated tokens, defined:

$$ u(x,y)=-\frac{1}{T}\sum_{t=1}^{T}\log(\widetilde p_t) $$

Here $\widetilde p_t$ is the generated token's probability after normalization.

The simulator sweeps a common deferral-rate parameter $r\in\{0,0.02,\ldots,1\}$. At each nonterminal position, it ranks uncertainty among the queries that actually reached that position and defers the highest-ranked fraction using the same rule. In the routed tier, this ranking is performed separately for each selected intervention's arrivals. At the initial position, deferral invokes the selector; after an intervention, deferral invokes $L$. The small-to-large baseline applies the identical uncertainty statistic and ranking rule but defers directly to $M_L$.

This procedure traces allocation curves under controlled deferral rates; it is an offline ranking evaluation, not a claim that a fixed numerical threshold has been calibrated for deployment. Random deferral and the notebook's learned correctness estimators can be reported as supplementary comparisons, while the primary architectural comparison holds the non-learned uncertainty rule fixed. The selector's action ranking and the gate's decision to continue are therefore evaluated as distinct components.

\subsection{Cost accounting}

For action $a$ on query $i$, we use separate input and output charges,
\begin{equation}
c_{i,a}=\alpha_{m(a)}\ell^{\mathrm{in}}_{i,a}
       +\beta_{m(a)}\ell^{\mathrm{out}}_{i,a},
\end{equation}
where $m(a)$ identifies the generation model. The notebook assigns relative rates $(\alpha_S,\beta_S)=(1,6)$ to baseline and intervention calls, and $(\alpha_L,\beta_L)=(9,54)$ to the large model. These are simulation rates, not measured Colab charges or asserted provider prices. Equal rates do not imply equal call costs: retrieval, demonstrations, and verification can change both prompt and response length.

For a routed cascade, let $d_{i,0}$ indicate entry into the intervention tier and $d_{i,1}$ indicate subsequent escalation. Its generation cost is
\begin{equation}
C_i=c_{i,0}+d_{i,0}\left(c_{i,a_i}+d_{i,1}c_{i,L}\right).
\end{equation}
The initial call is always charged. Verification incurs an additional prompt and completion charge, with its initial answer already accounted for by $c_{i,0}$. For a sequential baseline, all visited intervention costs are added analogously.

The supplied replay code measures output length by generated-token count but currently measures input length by whitespace-separated words. Its current spend is consequently an input-length proxy, not exact token billing. The intended per-token analysis should replace that input proxy with the model tokenizer's length for the complete serialized prompt, including demonstrations, retrieved text, and verification context. The generation-cost model excludes retrieval infrastructure and selector inference; these overheads should be reported separately or included in a sensitivity analysis before making end-to-end efficiency claims.
% AUTHOR NOTE: Resolve the input-length proxy before final cost figures, and specify whether router/retrieval overhead is measured, modelled, or excluded. The existing token sequences may permit reconstructing actual input token counts without rerunning generation.

\subsection{Evaluation aligned with the contributions}

\paragraph{Intervention complementarity.}
Let $E_0=\{i:q_{i,0}=0\}$ be baseline errors. For each intervention, define its recovery set $R_a=\{i\in E_0:q_{i,a}=1\}$ and regression set $G_a=\{i:q_{i,0}=1,q_{i,a}=0\}$. We measure recovery as $|R_a|/|E_0|$ and regressions relative to the initially correct set. Pairwise Jaccard overlap, $|R_a\cap R_b|/|R_a\cup R_b|$, and each action's uniquely recovered errors characterize complementarity; empty unions are reported as undefined. Applying the same definitions to $L$ distinguishes errors recovered by both mechanisms from those recovered only by interventions or only by the large model. We report these quantities for the full workload and by task family.

\paragraph{Sequential execution.}
The comparison includes the direct $S\!\rightarrow\!L$ cascade, each single-intervention cascade $S\!\rightarrow\!a\!\rightarrow\!L$, a fixed sequential traversal of all interventions before $L$, and the routed intervention cascade. Sequential replay pays for every visited call and uses the same deferral signal. It reuses the recorded intervention outputs: verification conditions on the baseline answer, and retrieval and few-shot generation use their own contexts. This evaluates sequential execution of alternative actions, not cumulative prompt refinement in which every action consumes its predecessor's response. Such history-dependent stacking would require separately generated trajectory outcomes.

\paragraph{Suggested evaluation extensions.}
To avoid attributing a poor fixed order to sequential execution in general, evaluate all six intervention orders, with any selected order chosen on validation data. Compare the learned selector with a random action selector and a validation-selected fixed action under the same gates. An oracle selector using recorded correctness provides a diagnostic ceiling for intervention selection, but is not deployable and should be labelled separately; at nonzero downstream deferral, the most accurate intermediate action is not necessarily an oracle for the entire cascade. A query-only selector ablation can isolate the value of observing the baseline answer.

\paragraph{Quality-cost comparison.}
For policy $\pi$, report mean terminal correctness $Q_\pi=N^{-1}\sum_i q_{i,\pi}$ against mean accumulated cost $\overline C_\pi=N^{-1}\sum_i C_{i,\pi}$, alongside intervention usage and large-model deferral. The principal comparison is the nondominated quality-cost frontier, rather than accuracy at matched deferral rates, since different paths incur different costs. Suggested summaries are quality at common budgets and cost at common quality targets. Any area-under-curve summary should use the same cost interval and normalization across policies. Paired, task-stratified bootstrap intervals and repeated selector seeds would quantify uncertainty; they should precede claims of statistically significant Pareto improvement. A separate cost-ratio sensitivity simulation can test how sequential and routed architectures behave as the available cost range narrows.
% AUTHOR NOTE: The extensions, bootstrap analysis, common-budget summaries, and cost-ratio simulation above are proposed evaluation choices, not claims that the supplied notebooks already implement or have completed them.





\section{Typefaces}



\section{Title Information}


\section{Authors and Affiliations}


\section{Rights Information}



\section{Sectioning Commands}


\subsection{Subsection}
\label{sec:subsection}

This is a subsection.

\subsubsection{Subsubsection}
\label{sec:subsubsection}

This is a subsubsection.

\paragraph{Paragraph}

This is a paragraph.

\subparagraph{Subparagraph}



\section{Tables}



\begin{table*}
  \caption{Some Typical Commands}
  \label{tab:commands}
  \begin{tabular}{ccl}
    \toprule
    Command &A Number & Comments\\
    \midrule
    \texttt{{\char'134}author} & 100& Author \\
    \texttt{{\char'134}table}& 300 & For tables\\
    \texttt{{\char'134}table*}& 400& For wider tables\\
    \bottomrule
  \end{tabular}
\end{table*}


\section{Math Equations}




\section{Citations and Bibliographies}




%%
%% The acknowledgments section is defined using the "acks" environment
%% (and NOT an unnumbered section). This ensures the proper
%% identification of the section in the article metadata, and the
%% consistent spelling of the heading.
\begin{acks}
ACKS
\end{acks}


\section*{Ethics and Privacy Statement}


%%
%% The next two lines define the bibliography style to be used, and
%% the bibliography file.
\bibliographystyle{ACM-Reference-Format}
\bibliography{references}


%%
%% If your work has an appendix, this is the place to put it.
\appendix
The configured selector is trained for three epochs with batch size eight, learning rate $2\times10^{-5}$, and a maximum input length of 512 tokens. The implementation uses Adafactor, cosine learning-rate scheduling, and gradient clipping at norm 1.0. The primary three-way selector uses text only, with no numerical confidence features and no separate keep class. Acceptance remains the responsibility of the uncertainty gate.
\section{Research Methods}





\end{document}
\endinput
%%
%% End of file `sigconf.tex'.
